\section{Results}
\label{sec:results}

The manuscript tables are generated from the raw per-seed JSON files by \texttt{experiments/reproduce.py} and \texttt{paper/make\_figures.py}. Table~\ref{tab:main} reports global IID baseline performance; Table~\ref{tab:algorithms} compares algorithms on the two-regime datasets; Table~\ref{tab:partitions} compares aggregation methods under partition shift; and Table~\ref{tab:clients} reports per-client performance. Entries are mean $\pm$ sample standard deviation; $n$ is shown for global rows. A paired exact sign-flip test compares methods on the same seed and partition (Table~\ref{tab:significance}). With five seeds, the minimum attainable two-sided exact $p$-value is $0.0625$; these tests therefore have limited resolution.

Centralized training has the lowest RMSE on all four IID datasets (\SuiteCentralIidRmseMin{}--\SuiteCentralIidRmseMax{}), while FedAvg improves over local-only on FD001, FD003 and FD004 but is nearly tied with local-only on FD002. On the FD002/FD004 regime partitions, FedProx is close to FedAvg. FedAdam with the fixed server learning rate used here degrades substantially (RMSE \SuiteFdTwoRegimeFedAdamRmse{} on FD002 and \SuiteFdFourRegimeFedAdamRmse{} on FD004); this is a negative result for this configuration, not evidence that all FedAdam settings fail. The paired suite's capped exponential objective lowers aggregate NASA score relative to the corresponding MSE baseline on both regime datasets, with higher RMSE. The compact CNN-LSTM performs worse than the 1D-CNN under this training budget. None of the five-seed paired tests reaches $p<0.05$, consistent with the limited resolution of an exact test at this sample size.

\begin{table*}[t]
\centering
\caption{Centralized, isolated local and FedAvg baselines on IID partitions. Values are mean $\pm$ sample standard deviation across five seeds. NASA scores are summed over engines and should only be compared within the same dataset.}
\label{tab:main}
\begin{tabular}{llcrr}\toprule
Dataset & Method & Privacy & RMSE (cycles) & NASA score \\\midrule
\multirow{4}{*}{FD001} & Constant baseline & n/a & $40.46\pm0.00$ & $17\,301.19\pm0.00$ \\
 & Isolated local (avg.) & local & $15.55\pm0.46$ & $464.43\pm74.67$ \\
 & FedAvg swarm (ours) & local & \textbf{$14.92\pm0.42$} & \textbf{$394.85\pm45.73$} \\
 & Centralized (upper bound) & exposed & $14.75\pm0.69$ & $425.36\pm94.56$ \\
\midrule
\multirow{4}{*}{FD002} & Constant baseline & n/a & $43.41\pm0.00$ & $59\,296.05\pm0.00$ \\
 & Isolated local (avg.) & local & $24.96\pm0.99$ & $10\,702.21\pm2\,404.11$ \\
 & FedAvg swarm (ours) & local & \textbf{$21.98\pm0.99$} & \textbf{$4\,877.92\pm2\,291.46$} \\
 & Centralized (upper bound) & exposed & $19.57\pm0.86$ & $2\,974.50\pm1\,405.23$ \\
\bottomrule\end{tabular}

\end{table*}

\begin{table*}[t]
\centering
\caption{Global algorithm comparison on regime-grouped FD002 and FD004. ``exp.'' denotes the paired suite's capped exponential NASA-score training objective, not \texttt{AsymmetricRULLoss}.}
\label{tab:algorithms}
\begin{tabular}{lllrrr}
\toprule
Dataset & Partition & Method & $n$ & RMSE & NASA score \\
\midrule
FD002 & regime & Centralized & 5 & $21.46\pm1.18$ & $5\,762.79\pm2\,135.75$ \\
FD002 & regime & Local only & 5 & $23.04\pm0.43$ & $5\,033.02\pm818.81$ \\
FD002 & regime & FedAvg & 5 & $23.50\pm0.63$ & $4\,425.37\pm1\,367.59$ \\
FD002 & regime & FedProx & 5 & $23.30\pm0.48$ & $4\,581.90\pm1\,331.95$ \\
FD002 & regime & FedAdam & 5 & $78.23\pm23.17$ & $3.29\times10^{8}\pm6.99\times10^{8}$ \\
FD002 & regime & FedAvg + exp. & 5 & $27.44\pm0.69$ & $4\,104.37\pm534.95$ \\
FD002 & regime & Central + exp. & 5 & $25.62\pm1.11$ & $4\,009.00\pm422.46$ \\
FD002 & regime & CNN-LSTM & 5 & $43.14\pm0.02$ & $25\,313.76\pm294.55$ \\
FD004 & regime & Centralized & 5 & $23.41\pm0.49$ & $4\,329.01\pm1\,650.39$ \\
FD004 & regime & Local only & 5 & $26.85\pm0.80$ & $8\,006.51\pm2\,168.07$ \\
FD004 & regime & FedAvg & 5 & $26.78\pm0.20$ & $7\,654.82\pm693.72$ \\
FD004 & regime & FedProx & 5 & $26.88\pm0.81$ & $7\,638.47\pm1\,490.37$ \\
FD004 & regime & FedAdam & 5 & $56.62\pm16.01$ & $9.00\times10^{6}\pm1.84\times10^{7}$ \\
FD004 & regime & FedAvg + exp. & 5 & $29.75\pm0.36$ & $6\,513.13\pm332.91$ \\
FD004 & regime & Central + exp. & 5 & $29.95\pm0.82$ & $7\,243.98\pm3\,590.53$ \\
FD004 & regime & CNN-LSTM & 5 & $43.00\pm0.01$ & $43\,970.10\pm652.10$ \\
\bottomrule
\end{tabular}

\end{table*}

\begin{table*}[t]
\centering
\caption{FedAvg, FedProx and FedAdam across IID, regime-grouped and engine-count-imbalanced partitions on FD002 and FD004.}
\label{tab:partitions}
\begin{tabular}{lllrrr}
\toprule
Dataset & Partition & Method & $n$ & RMSE & NASA score \\
\midrule
FD002 & iid & fedavg & 5 & $23.33\pm0.24$ & $4\,179.66\pm680.06$ \\
FD002 & regime & fedavg & 5 & $23.50\pm0.63$ & $4\,425.37\pm1\,367.59$ \\
FD002 & regime & fedprox & 5 & $23.30\pm0.48$ & $4\,581.90\pm1\,331.95$ \\
FD002 & regime & fedadam & 5 & $78.23\pm23.17$ & $328\,950\,600.23\pm699\,140\,126.97$ \\
FD002 & engine\_imbalance & fedavg & 5 & $22.36\pm0.39$ & $4\,073.82\pm970.12$ \\
FD002 & engine\_imbalance & fedprox & 5 & $22.30\pm0.82$ & $4\,170.31\pm1\,185.30$ \\
FD002 & engine\_imbalance & fedadam & 5 & $59.37\pm16.21$ & $10\,521\,602.78\pm20\,759\,533.05$ \\
FD004 & iid & fedavg & 5 & $27.07\pm0.56$ & $8\,059.20\pm2\,868.94$ \\
FD004 & regime & fedavg & 5 & $26.78\pm0.20$ & $7\,654.82\pm693.72$ \\
FD004 & regime & fedprox & 5 & $26.88\pm0.81$ & $7\,638.47\pm1\,490.37$ \\
FD004 & regime & fedadam & 5 & $56.62\pm16.01$ & $8\,996\,132.47\pm18\,447\,615.59$ \\
FD004 & engine\_imbalance & fedavg & 5 & $26.55\pm0.73$ & $8\,205.18\pm4\,119.58$ \\
FD004 & engine\_imbalance & fedprox & 5 & $26.33\pm0.77$ & $7\,748.58\pm3\,981.96$ \\
FD004 & engine\_imbalance & fedadam & 5 & $49.66\pm12.60$ & $2\,137\,084.85\pm4\,493\,691.96$ \\
\bottomrule
\end{tabular}

\end{table*}

\begin{table*}[t]
\centering
\caption{Per-client RMSE and NASA score on regime-grouped FD002 and FD004 (mean $\pm$ sample standard deviation over five paired seeds). Client labels denote the two operating-regime groups; ``exp.'' denotes the capped exponential objective.}
\label{tab:clients}
\begin{tabular}{llllrr}
\toprule
Dataset & Partition & Method & Client & RMSE & NASA score \\
\midrule
FD002 & regime & centralized & 1 & $21.78\pm1.12$ & $1\,824.32\pm867.07$ \\
FD002 & regime & centralized & 2 & $21.21\pm1.23$ & $3\,938.47\pm1\,313.34$ \\
FD002 & regime & local\_only & 1 & $23.51\pm0.43$ & $1\,811.16\pm587.92$ \\
FD002 & regime & local\_only & 2 & $22.67\pm0.59$ & $3\,221.86\pm1\,129.69$ \\
FD002 & regime & fedavg & 1 & $23.61\pm0.67$ & $1\,594.57\pm502.96$ \\
FD002 & regime & fedavg & 2 & $23.41\pm0.60$ & $2\,830.80\pm870.10$ \\
FD002 & regime & fedprox & 1 & $23.35\pm0.58$ & $1\,589.67\pm474.67$ \\
FD002 & regime & fedprox & 2 & $23.26\pm0.41$ & $2\,992.23\pm862.77$ \\
FD002 & regime & fedadam & 1 & $77.42\pm23.30$ & $126\,861\,499.80\pm269\,199\,499.94$ \\
FD002 & regime & fedadam & 2 & $78.84\pm23.07$ & $202\,089\,088.31\pm429\,940\,704.72$ \\
FD002 & regime & federated\_asymmetric & 1 & $26.85\pm0.69$ & $1\,559.59\pm141.17$ \\
FD002 & regime & federated\_asymmetric & 2 & $27.88\pm0.74$ & $2\,544.79\pm408.81$ \\
FD002 & regime & asymmetric & 1 & $25.18\pm1.16$ & $1\,479.20\pm74.43$ \\
FD002 & regime & asymmetric & 2 & $25.96\pm1.08$ & $2\,529.80\pm364.81$ \\
FD002 & regime & cnn\_lstm & 1 & $41.43\pm0.02$ & $10\,189.80\pm121.04$ \\
FD002 & regime & cnn\_lstm & 2 & $44.40\pm0.02$ & $15\,123.95\pm173.50$ \\
FD004 & regime & centralized & 1 & $20.61\pm0.55$ & $825.25\pm102.61$ \\
FD004 & regime & centralized & 2 & $25.18\pm0.68$ & $3\,503.76\pm1\,560.42$ \\
FD004 & regime & local\_only & 1 & $23.85\pm0.68$ & $1\,160.84\pm191.71$ \\
FD004 & regime & local\_only & 2 & $28.75\pm1.04$ & $6\,845.68\pm2\,061.46$ \\
FD004 & regime & fedavg & 1 & $23.55\pm0.29$ & $1\,186.92\pm64.29$ \\
FD004 & regime & fedavg & 2 & $28.81\pm0.25$ & $6\,467.90\pm642.30$ \\
FD004 & regime & fedprox & 1 & $23.73\pm1.01$ & $1\,192.59\pm184.98$ \\
FD004 & regime & fedprox & 2 & $28.88\pm0.74$ & $6\,445.88\pm1\,328.91$ \\
FD004 & regime & fedadam & 1 & $58.99\pm16.31$ & $4\,860\,573.40\pm9\,971\,441.12$ \\
FD004 & regime & fedadam & 2 & $54.89\pm15.80$ & $4\,135\,558.69\pm8\,476\,173.80$ \\
FD004 & regime & federated\_asymmetric & 1 & $28.12\pm0.40$ & $1\,700.15\pm77.67$ \\
FD004 & regime & federated\_asymmetric & 2 & $30.83\pm0.36$ & $4\,812.98\pm260.78$ \\
FD004 & regime & asymmetric & 1 & $28.19\pm0.94$ & $1\,820.55\pm628.55$ \\
FD004 & regime & asymmetric & 2 & $31.12\pm0.77$ & $5\,423.42\pm2\,966.02$ \\
FD004 & regime & cnn\_lstm & 1 & $44.24\pm0.01$ & $23\,522.60\pm361.45$ \\
FD004 & regime & cnn\_lstm & 2 & $42.12\pm0.00$ & $20\,447.51\pm290.64$ \\
\bottomrule
\end{tabular}

\end{table*}

\begin{table*}[t]
\centering
\caption{Paired exact sign-flip tests for baseline and method comparisons, by dataset, partition and metric.}
\label{tab:significance}
\begin{tabular}{llllr}
\toprule
Dataset & Partition & Comparison & Metric & paired $p$ \\
\midrule
FD001 & iid & fedavg\_vs\_centralized & rmse & 0.0625 \\
FD001 & iid & fedavg\_vs\_centralized & nasa & 0.0625 \\
FD001 & iid & fedavg\_vs\_local\_only & rmse & 0.1875 \\
FD001 & iid & fedavg\_vs\_local\_only & nasa & 0.0625 \\
FD002 & regime & fedavg\_vs\_centralized & rmse & 0.1250 \\
FD002 & regime & fedavg\_vs\_centralized & nasa & 0.3750 \\
FD002 & regime & fedavg\_vs\_local\_only & rmse & 0.2500 \\
FD002 & regime & fedavg\_vs\_local\_only & nasa & 0.5000 \\
FD002 & regime & fedavg\_vs\_fedprox & rmse & 0.1250 \\
FD002 & regime & fedavg\_vs\_fedprox & nasa & 0.0625 \\
FD002 & regime & fedavg\_vs\_fedadam & rmse & 0.0625 \\
FD002 & regime & fedavg\_vs\_fedadam & nasa & 0.0625 \\
FD002 & regime & fedavg\_vs\_federated\_asymmetric & rmse & 0.0625 \\
FD002 & regime & fedavg\_vs\_federated\_asymmetric & nasa & 0.5625 \\
FD002 & regime & centralized\_vs\_asymmetric & rmse & 0.0625 \\
FD002 & regime & centralized\_vs\_asymmetric & nasa & 0.1875 \\
FD002 & regime & centralized\_vs\_cnn\_lstm & rmse & 0.0625 \\
FD002 & regime & centralized\_vs\_cnn\_lstm & nasa & 0.0625 \\
FD004 & regime & fedavg\_vs\_centralized & rmse & 0.0625 \\
FD004 & regime & fedavg\_vs\_centralized & nasa & 0.0625 \\
FD004 & regime & fedavg\_vs\_local\_only & rmse & 0.8125 \\
FD004 & regime & fedavg\_vs\_local\_only & nasa & 0.7500 \\
FD004 & regime & fedavg\_vs\_fedprox & rmse & 1.0000 \\
FD004 & regime & fedavg\_vs\_fedprox & nasa & 1.0000 \\
FD004 & regime & fedavg\_vs\_fedadam & rmse & 0.0625 \\
FD004 & regime & fedavg\_vs\_fedadam & nasa & 0.0625 \\
FD004 & regime & fedavg\_vs\_federated\_asymmetric & rmse & 0.0625 \\
FD004 & regime & fedavg\_vs\_federated\_asymmetric & nasa & 0.0625 \\
FD004 & regime & centralized\_vs\_asymmetric & rmse & 0.0625 \\
FD004 & regime & centralized\_vs\_asymmetric & nasa & 0.3125 \\
FD004 & regime & centralized\_vs\_cnn\_lstm & rmse & 0.0625 \\
FD004 & regime & centralized\_vs\_cnn\_lstm & nasa & 0.0625 \\
\bottomrule
\end{tabular}

\end{table*}

\subsection{Asymmetric risk-aware training}
Table~\ref{tab:asymmetric} compares symmetric \texttt{mse} against quadratic \texttt{AsymmetricRULLoss} ($\alpha_{\text{late}}{=}2.5$, $\alpha_{\text{early}}{=}0.5$) under the separate controlled five-seed protocol in \texttt{benchmark\_asymmetric\_loss.py}. Because this protocol uses different seeds, preprocessing and batch size, its absolute values should not be compared with Tables~\ref{tab:main}--\ref{tab:clients}.

\begin{table*}[t]
\centering
\caption{Controlled five-seed comparison of symmetric MSE vs.\ \texttt{AsymmetricRULLoss} ($\alpha_{\text{late}}{=}2.5$, $\alpha_{\text{early}}{=}0.5$; seeds 42, 123, 456, 789, 1024; mean $\pm$ std). Source: \texttt{results/asymmetric\_loss\_comparison.json}. Bold marks the better value per setup and dataset.}
\label{tab:asymmetric}
\begin{tabular}{llcccrr}\toprule
Dataset & Setup & Loss & RMSE (cycles) & NASA score & $\Delta$ NASA & Late rate (\%) \\\midrule
\multirow{6}{*}{FD001} & \multirow{2}{*}{Isolated local} & \texttt{mse} & $17.58\pm1.01$ & $722.9\pm120.5$ & --- & $56.2\pm6.6$ \\
 &  & \texttt{asymmetric} & \textbf{$17.15\pm1.50$} & \textbf{$435.0\pm85.6$} & $-39.8\%$ & \textbf{$33.8\pm4.5$} \\
\cmidrule(lr){2-7}
 & \multirow{2}{*}{FedAvg (5r $\times$ 2e)} & \texttt{mse} & $16.98\pm1.31$ & $690.8\pm173.4$ & --- & $63.8\pm7.1$ \\
 &  & \texttt{asymmetric} & \textbf{$16.25\pm1.23$} & \textbf{$369.1\pm43.5$} & $-46.6\%$ & \textbf{$34.2\pm8.0$} \\
\cmidrule(lr){2-7}
 & \multirow{2}{*}{Centralized (10e)} & \texttt{mse} & \textbf{$15.16\pm0.86$} & $450.4\pm91.8$ & --- & $56.6\pm3.3$ \\
 &  & \texttt{asymmetric} & $15.80\pm0.97$ & \textbf{$338.7\pm52.0$} & $-24.8\%$ & \textbf{$30.8\pm5.2$} \\
\midrule
\multirow{6}{*}{FD002} & \multirow{2}{*}{Isolated local} & \texttt{mse} & \textbf{$23.73\pm1.22$} & $7\,610.8\pm3\,611.4$ & --- & $50.3\pm7.8$ \\
 &  & \texttt{asymmetric} & $26.30\pm2.08$ & \textbf{$3\,932.2\pm829.7$} & $-48.3\%$ & \textbf{$29.5\pm6.0$} \\
\cmidrule(lr){2-7}
 & \multirow{2}{*}{FedAvg (5r $\times$ 2e)} & \texttt{mse} & \textbf{$22.33\pm1.46$} & $5\,499.1\pm2\,962.8$ & --- & $48.9\pm8.7$ \\
 &  & \texttt{asymmetric} & $26.07\pm1.69$ & \textbf{$3\,412.6\pm429.8$} & $-37.9\%$ & \textbf{$25.3\pm6.3$} \\
\cmidrule(lr){2-7}
 & \multirow{2}{*}{Centralized (10e)} & \texttt{mse} & \textbf{$20.17\pm0.31$} & \textbf{$2\,769.5\pm494.6$} & --- & $40.8\pm5.0$ \\
 &  & \texttt{asymmetric} & $25.79\pm2.74$ & $3\,630.2\pm1\,185.2$ & $+31.1\%$ & \textbf{$20.5\pm6.1$} \\
\bottomrule\end{tabular}

\end{table*}

Across all six dataset--setup pairs, \texttt{AsymmetricRULLoss} reduces the late-prediction rate ($\hat{y} > y$) by \AsymLateReductionMin{}\%--\AsymLateReductionMax{}\%, from \AsymMseLateMin{}\%--\AsymMseLateMax{}\% under MSE to \AsymLateMin{}\%--\AsymLateMax{}\%. In decentralized settings, this shift reduces both the mean and seed-to-seed variance of the NASA score: on FD001, FedAvg drops by \AsymFdOneFedNasaDelta{} (\AsymFdOneFedMseNasa{}$\pm$\AsymFdOneFedMseNasaStd{} to \AsymFdOneFedAsyNasa{}$\pm$\AsymFdOneFedAsyNasaStd{}) and Isolated Local by \AsymFdOneLocalNasaDelta{}; on FD002, FedAvg drops by \AsymFdTwoFedNasaDelta{} (\AsymFdTwoFedMseNasa{}$\pm$\AsymFdTwoFedMseNasaStd{} to \AsymFdTwoFedAsyNasa{}$\pm$\AsymFdTwoFedAsyNasaStd{}, with standard deviation falling by a factor of \AsymFdTwoFedNasaStdRatio{}) and Isolated Local by \AsymFdTwoLocalNasaDelta{}.

The results also expose a dataset-dependent safety--accuracy trade-off. On FD001, \texttt{AsymmetricRULLoss} improves both NASA score and RMSE for Isolated Local (\AsymFdOneLocalMseRmse{} to \AsymFdOneLocalAsyRmse{}) and FedAvg (\AsymFdOneFedMseRmse{} to \AsymFdOneFedAsyRmse{}), with only a slight RMSE increase for Centralized (\AsymFdOneCentMseRmse{} to \AsymFdOneCentAsyRmse{}). On FD002, however, the conservative early-prediction bias increases symmetric RMSE across all setups (\AsymFdTwoLocalMseRmse{} to \AsymFdTwoLocalAsyRmse{} Local; \AsymFdTwoFedMseRmse{} to \AsymFdTwoFedAsyRmse{} FedAvg; \AsymFdTwoCentMseRmse{} to \AsymFdTwoCentAsyRmse{} Centralized) and raises the Centralized NASA score (\AsymFdTwoCentNasaDelta{}, driven by conservative early-prediction outliers such as seed 1024).

\subsection{Edge cost}
Per-round CPU training time, peak process resident memory, and model-tensor payload bytes are recorded by \texttt{experiments/profile\_edge.py}. The committed one-round, one-epoch FD002 smoke profile on Apple Silicon records \EdgeSmokeTrainSeconds{}~s, \EdgeSmokePeakRssMiB{}~MiB peak process RSS and \EdgeSmokePayloadBytes{} tensor bytes for one download plus one upload. It is not evidence of Raspberry Pi performance; no Pi profile completed because the remote hosts were inaccessible. Payload excludes gRPC framing and therefore is not measured network wire traffic.

\subsection{Result provenance}
Raw metrics and training configuration are stored in \texttt{results/reproducible/}. The older FD001/FD002 artifacts use a different pipeline and are not included in these tables.
